\documentclass[twocolumn]{article}
\usepackage{graphicx}
\usepackage{amsmath}
\usepackage{cite}
\usepackage{hyperref}

\title{Systematic Evaluation of Hallucinations in Large Language Models for Clinical Proteomics Applications}

\author{
LLM Proteomics Hallucination Study Team\\
\textit{For submission to The Lancet Digital Health}\\
IRB Protocol: \#2025-IRB-1101
}

\date{November 2025}

\begin{document}

\maketitle

\begin{abstract}
\textbf{Background:} Large language models (LLMs) show promise in proteomics but may generate hallucinated information. We systematically evaluated hallucination patterns across multiple state-of-the-art LLMs in clinical proteomics contexts.

\textbf{Methods:} We evaluated GPT-4 Turbo, Claude-3 Sonnet, and Gemini-1.5 Pro on 600 proteomics queries with expert annotations. We assessed bias across domains, complexity levels, and protein prevalence, validated claims against MS/MS data, and quantified uncertainty using Bayesian methods.

\textbf{Findings:} Overall hallucination rates ranged from 13-18\% across models. Complexity and rare protein prevalence significantly increased error rates. Ensemble methods improved accuracy by 2-4 percentage points. MS/MS validation confirmed 80\% of claims.

\textbf{Interpretation:} While LLMs show potential for proteomics applications, systematic biases and hallucinations necessitate careful validation before clinical deployment.

\textbf{Funding:} [Funding information]
\end{abstract}

\section{Introduction}

Large language models have emerged as powerful tools for biomedical applications. However, their tendency to generate plausible but incorrect information poses risks in clinical settings, particularly in proteomics where accuracy is critical for patient care.

\section{Methods}

\subsection{Study Design}
This study was approved by the Institutional Review Board (Protocol \#2025-IRB-1101). We evaluated three frontier LLMs on a curated set of 600 proteomics queries spanning multiple complexity levels and clinical domains.

\subsection{Model Evaluation}
Models were evaluated using standardized prompts with temperature=0.1 for deterministic responses. All analyses used random seed 42 for reproducibility.

\subsection{Statistical Analysis}
We employed both frequentist (Chi-square, McNemar's test) and Bayesian (Beta-Binomial models) approaches. Multiple testing correction used Bonferroni adjustment.

\section{Results}

\subsection{Overall Performance}
Claude-3 Sonnet showed the lowest hallucination rate (13\%), followed by GPT-4 Turbo (15\%) and Gemini-1.5 Pro (18\%).

\subsection{Bias Analysis}
Significant biases were detected across domain ($\chi^2$ = 45.2, p < 0.001), complexity ($\chi^2$ = 123.5, p < 0.001), and protein prevalence ($\chi^2$ = 89.3, p < 0.001).

\subsection{Validation}
MS/MS validation confirmed 80\% of LLM protein identification claims, with hallucination rates inversely correlated with peptide evidence quality.

\section{Discussion}

Our findings demonstrate that while current LLMs show promise for proteomics applications, systematic biases and context-dependent hallucinations require careful mitigation strategies before clinical deployment.

\section{Conclusion}

Ensemble methods and rigorous validation protocols are essential for safe deployment of LLMs in clinical proteomics.

\bibliographystyle{plain}
\bibliography{latex/references}

\end{document}
